\documentclass[letterpaper]{exam}
\usepackage[utf8]{inputenc}
\usepackage[spanish]{babel}
\usepackage[voffset=0.5in,margin=0.8in]{geometry}
\usepackage{amsmath,amssymb}
\usepackage{amsthm}
\usepackage{amsfonts}
\usepackage{mathptmx}
\usepackage{stmaryrd}
\usepackage{mathpartir}
\usepackage{mathtools}
\usepackage{listings}
\usepackage{fancyvrb}
\usepackage{multicol}
\usepackage[hidelinks]{hyperref}
\usepackage[dvipsnames,table]{xcolor}
\usepackage{graphicx}
\usepackage{booktabs}
\usepackage{makecell}
\usepackage{enumitem}
\usepackage{logicproof}
\usepackage{tikz}
\usepackage{fitch}
% \usepackage{emoji}
\usepackage{multirow}
\definecolor{light-gray}{gray}{0.9}

\graphicspath{{../images/}}

\newcommand{\fcal}{\mathcal{F}}
\newcommand{\zmat}{\mathbb{Z}}
\newcommand{\true}{\textbf{\textsf{T}}}
\newcommand{\false}{\textbf{\textsf{F}}}
\newcommand{\vtwo}{45}
\newcommand{\replaces}{\,\textbf{::}\,}

\newcommand{\boxd}[2]{
\fbox{%
\parbox{#1\textwidth}{\(#2\)}%
}
}

\newcommand{\class}{Curso}
\newcommand{\term}{Escuela de Ciencias Aplicadas e Ingeniería}
\newcommand{\examnum}{Parcial}
\newcommand{\examdate}{Fecha}
\newcommand{\prof}{Profesor}
\newcommand{\name}{\textbf{Nombre}:\enspace\makebox[2in]{\hrulefill}}
\newcommand{\timelimit}{1 hora y 20 minutos}
\renewcommand{\solutiontitle}{\noindent\textbf{Solución:}\par\noindent}

\firstpageheader{\includegraphics[scale=0.3]{logo-eafit.png}}{}{ \term\\\class\, - \examnum\\\prof\\\examdate}

\runningheader{\oddeven{\name}{\class}}{Pág. \thepage\ de \totalnumpages}{\examdate}
\runningheadrule{}

% Configuración de la tabla de calificación
\pointpoints{punto}{puntos}
\bonuspointpoints{bono}{bono}
\chpword{Puntos}
\vpword{Puntos}
\vtword{Total}
\htword{Total}
\vsword{Resultado}
\hsword{Resultado:}
\vqword{Problema}
\hqword{Pregunta}
\chqword{Pregunta}
\chbpword{Bono}
\chsword{Puntaje}
\marginbonuspointname{B}
\pointsinmargin{}
\boxedpoints{}
\setlength\fillinlinelength{1.5in}

\begin{document}
\noprintanswers{}

\noindent
\begin{tabular*}{\textwidth}{l @{\extracolsep{\fill}} l @{\extracolsep{6pt}}}
 \textbf{Nombre completo} & \textbf{Código estudiante o documento de identidad} \\
 & \\
 \makebox[3in]{\hrulefill} & \makebox[3in]{\hrulefill}
% \textbf{Tiempo: \timelimit} & Profesor: & \makebox[2.5in]{\emph{Dr. John Doe}}
\end{tabular*}
% \rule[2ex]{\textwidth}{2pt}

% Bloque de instrucciones
\begin{center}
\fbox{%
\parbox{0.98\textwidth}{
Este parcial contiene \numquestions\, preguntas que corresponden en total a \numpoints\, puntos y \numbonuspoints\, puntos de bono. El valor de cada pregunta se encuentra en un cuadro en el margen izquierdo de cada hoja. El parcial debe resolverse y entregarse individualmente. No está permitido hablar o pedir cosas prestadas de sus compañeros, tampoco el uso de libros, apuntes o aparatos electrónicos. El tiempo disponible para resolver el parcial es \timelimit\, y no se reciben hojas adicionales a las del enunciado.
}%
}
\end{center}

% Tabla de calificaciones
\begin{center}
\addpoints{}
% Puede presentar la tabla de calificación en orientación vertical [v] u horizontal [h]
\combinedgradetable[h][questions]
% \gradetable[h][questions]

\footnotesize{Tabla de calificación para uso exclusivo del profesor.}
\end{center}

\begin{table}[h]
  \centering
  \begin{tabular}{ llll }
  \toprule
  \multicolumn{4}{c}{\textbf{Reglas de inferencia}}\\
  \midrule
  Modus Ponens (MP) & \inferrule{p \to q \\ p}{q} & Dilema destructivo (DD) & \inferrule{(p \to q) \land(r \to s) \\ \neg q \lor \neg s}{\neg p \lor\neg r} \\\\
  Modus Tollens (MT) & \inferrule{p \to q \\ \neg q}{\neg p} & Simplificación (S) & \inferrule{p \land q}{p} \\\\
  Silogismo hipotético (SH) & \inferrule{p \to q \\ q \to r}{p \to r} & Conjunción (C) & \inferrule{p \\ q}{p \land q} \\\\
  Silogismo disyuntivo (SD) & \inferrule{p \lor q \\ \neg p}{q} & Adición (A) & \inferrule{p}{p \lor q} \\\\
  Dilema constructivo (DC) & \inferrule{(p \to q) \land(r \to s) \\ p \lor r}{q \lor s} && \\
  \bottomrule
  \end{tabular}
\end{table}

\begin{table}[h]
  \centering
  \begin{tabular}{ llll }
  \toprule
  \multicolumn{4}{c}{\textbf{Reglas de reemplazo}}\\
  \midrule
  Leyes de De Morgan (DM) & \thead{\(\neg (p \lor q) \replaces (\neg p \land \neg q)\)\\ \(\neg (p \land q) \replaces (\neg p \lor \neg q)\)} & Transposición (Trans.) & \((p \to q) \replaces (\neg q \to \neg p)\)\\\\
  Conmutatividad (Conm.) & \thead{\((p \land q) \replaces (q \land p)\)\\\((p \lor q) \replaces (q \lor p)\)} & Implicación material (IM) & \( (p \to q) \replaces (\neg p \lor q)\) \\\\
  Asociatividad (Asoc.) & \thead{\((p \land (q \land r)) \replaces ((p \land q) \land r)\)\\\((p \lor (q \lor r)) \replaces ((p \lor q) \lor r) \)} & Equivalencia material (EM) & \thead{\((p \equiv q) \replaces ((p \land q) \lor (\neg p \land \neg q ))\)\\\((p \equiv q) \replaces ((p \to q) \land (q \to p))\)}\\\\
  Distributividad (Dist.) & \thead{\((p \land (q \lor r)) \replaces ((p \land q) \lor (p \land r))\)\\\((p \lor (q \land r)) \replaces ((p \lor q) \land (p \lor r))\)} & Exportación (Exp.) & \( ((p \land q) \to r) \replaces (p \to (q \to r))\) \\\\
  Doble Negación (DN) & \(p \replaces \neg \neg p\) & Tautología (Taut.) & \thead{\(p \replaces (p \lor p)\)\\\(p \replaces (p \land p)\)}\\
  \bottomrule
  \end{tabular}
\end{table}
\newpage
\begin{questions}
  \question Pregunta 1.
    \boxd{0.4}{
    Cuadro-para-texto-matematico
    }
    \begin{parts}
        \part[15] \emph{Espacio para solución usando líneas punteadas}
        \fillwithdottedlines{\stretch{1}}
        \part[15] \emph{Espacio para solución:}
        \makeemptybox{\stretch{1}}
    \end{parts}
    \newpage
    \question Pregunta 2. Diferente espacio para solución, dejando simplemente espacio en blanco.
    \begin{parts}
        \part[5] Parte 1
        \vspace{\stretch{1}}
        \part[5] Parte 2.
        \vspace{\stretch{2}}
        \bonuspart[7] Parte bono.
        \vspace{\stretch{3}}
    \end{parts}
    \bonusquestion Pregunta bono.
    \fillwithdottedlines{\stretch{2}}
\end{questions}  
\end{document}
